\documentclass[11pt, a4paper]{article}
\usepackage[utf8]{inputenc}
\usepackage[ngerman]{babel}
\usepackage[T1]{fontenc}
\usepackage{amsmath, amssymb, amsfonts}
\usepackage{geometry}
\geometry{a4paper, left=2.5cm, right=2.5cm, top=3cm, bottom=3cm}
\usepackage{hyperref}
\usepackage{titlesec}
\usepackage{enumitem}
\usepackage{setspace}
\usepackage{booktabs}
\setstretch{1.15}
\titleformat{\section}{\large\bfseries\sffamily}{\thesection}{1em}{}
\newcommand{\mvec}[1]{\boldsymbol{#1}}
\newcommand{\n}{\mvec{n}}
\newcommand{\clc}{\mathcal{C}\ell_{3}(\mathbb{C})}

\begin{document}
\begin{center}
  \sffamily
  {\Large\textbf{Clifford-BEM: Ausarbeitung zu AP 3}}\\[0.3cm]
  {\large Teil I: $\mathcal{H}$-Matrix mit Adaptiver Kreuzapproximation;\\ kombinierter Löser am Würfel}\\[0.3cm]
  {\small Arbeitsstand, Version 0.5}
\end{center}

\section*{Überblick}
Grundlage ist der Galerkin-Prototyp aus AP~2 mit stückweise konstanten
Multivektor-Dichten auf ebenen Dreiecken. Der Randoperator $E_k$ wird hierarchisch
komprimiert (\texttt{hmat.py}), und die Formulierung $T_1$ wird mit GMRES über das
komprimierte Matrix-Vektor-Produkt gelöst (\texttt{scatter\_h.py}).

\paragraph{Ergebnisse.}
\begin{itemize}[leftmargin=*]
  \item Der Fehler des Matrix-Vektor-Produkts folgt der Toleranz: $\approx0{,}1\,\varepsilon$.
  \item Bei 3920 Dreiecken (31\,360 Unbekannte) braucht der Operator 270\,MB statt 15\,GB
    für die dichte Matrix.
  \item Die Streurechnung auf der Kugel ist jetzt bis 23\,040 Unbekannte möglich. Für Glas
    ergibt die Extrapolation $Q_{\mathrm{ext}}=0{,}21511$ gegenüber Mie $0{,}21510$.
  \item Die gemeinsame Kreuzapproximation aller Kernkomponenten spart gegenüber der
    komponentenweisen 10 bis 19\,\% Speicher bei gleicher Genauigkeit (H3, erster Teil).
  \item Nahezu lineare Komplexität (H3, zweiter Teil): Im Python-Bereich (bis 3920
    Dreiecke) wuchs der Speicher noch wie $N^{1{,}39}$. Mit dem C++-Kern bis 20\,480 Dreiecke
    zeigt sich $\mathcal{O}(N\log^2N)$ (Abschnitt~\ref{sec:cpp}).
  \item Am Würfel arbeiten Kompression und Kanten-/Eck-Vorkonditionierung zusammen: Die
    GMRES-Iterationen sind von der Kantenauflösung unabhängig (Abschnitt~\ref{sec:kombi}).
\end{itemize}

\section{Konzept}

\paragraph{Was komprimiert wird.} Für ein Dreieckspaar ist der Galerkin-Block
\[
  E_{ij} = -\frac{2}{\sqrt{|\tau_i||\tau_j|}}\,L\big((s_{ij}+\mvec{v}_{ij})\,\n_j\big),\qquad
  (s_{ij},\mvec{v}_{ij}) = \int_{\tau_i}\!\int_{\tau_j}\mvec\Phi_k(x-y)\,dS_y\,dS_x .
\]
Gespeichert werden nur die vier Kernkomponenten $K_c(i,j)$, also $s$ und die drei Komponenten
von $\mvec v$, nicht die $8\times8$-Blöcke. Weil $E_{ij}$ linear in $K_c$ ist, lässt sich das
Matrix-Vektor-Produkt schreiben als
\[
  (Ex)_i = -\frac{2}{\sqrt{|\tau_i|}}\sum_{c=0}^{3}\sum_j K_c(i,j)\,z^c_j,\qquad
  z^c_j = \frac{L(\mvec{e}_c\n_j)\,x_j}{\sqrt{|\tau_j|}},\quad \mvec{e}_0=1 .
\]
Das spart gegenüber den Blöcken den Faktor 16, schon bevor komprimiert wird.

\paragraph{Clusterbaum und Zulässigkeit.} Der Clusterbaum entsteht durch Halbierung
entlang der längsten Achse der Schwerpunkte, mit höchstens 32 Dreiecken je Blatt. Ein
Blockpaar $(\tau,\sigma)$ ist zulässig, wenn
$\min(\operatorname{diam}\tau,\operatorname{diam}\sigma)\le\eta\operatorname{dist}(\tau,\sigma)$
mit $\eta=1$ gilt. Außerdem muss $\operatorname{dist}>3h_{\max}$ sein, damit alle Paare im
Fernfeld der Gauß-Regel liegen, und $k\operatorname{diam}\le20$. Unzulässige Blätter werden
dicht mit der Nahfeldbehandlung aus AP~2 assembliert; diese ist jetzt über die
Nachbardreiecke vektorisiert.

\paragraph{Kreuzapproximation.} Verwendet wird ACA mit partieller Pivotisierung und
Abbruch bei $\|u_k\|\|v_k\|\le\varepsilon\|S_k\|_F$, danach eine QR/SVD-Nachkompression.
Bei komplexen Faktoren muss es in der Nachkompression $V=Q_VZ^{T}$ heißen und nicht
$Q_VZ^{H}$; die konjugierte Variante erzeugte Blockfehler von 20 bis 75\,\%.

\paragraph{Zwei Varianten.}
\begin{itemize}[leftmargin=*]
  \item \emph{komponentenweise:} vier getrennte ACAs für die vier Kernkomponenten. Speicher
    $\sum_c r_c(|\tau|+|\sigma|)$.
  \item \emph{gemeinsam:} eine ACA auf der gestapelten Matrix $|\tau|\times4|\sigma|$. Die
    Zeilenfaktoren und Pivots werden also über die Grade geteilt. Speicher
    $r(|\tau|+4|\sigma|)$.
\end{itemize}
Eine echte Block-ACA mit $8\times8$-Pivotblöcken ist noch nicht umgesetzt.

\section{Genauigkeit und Variantenvergleich}
Kugel mit 1280 Dreiecken, $k=1{,}5$. Der Fehler ist der relative Fehler des
Matrix-Vektor-Produkts gegenüber 40 exakt assemblierten Zeilen
(\texttt{test\_hmat\_large.py}):
\begin{center}\small
\begin{tabular}{lcccccc}
\toprule
& \multicolumn{3}{c}{gemeinsam} & \multicolumn{3}{c}{komponentenweise}\\
$\varepsilon$ & Fehler & Niedrigrang & Aufbau & Fehler & Niedrigrang & Aufbau\\
\midrule
$10^{-3}$ & $9{,}4\cdot10^{-5}$ & 18,8\,MB & 35\,s & $9{,}5\cdot10^{-5}$ & 20,9\,MB & 53\,s\\
$10^{-4}$ & $1{,}1\cdot10^{-5}$ & 25,1\,MB & 37\,s & $9{,}8\cdot10^{-6}$ & 31,1\,MB & 62\,s\\
$10^{-5}$ & $1{,}1\cdot10^{-6}$ & 35,0\,MB & 40\,s & $1{,}0\cdot10^{-6}$ & 41,8\,MB & 72\,s\\
\bottomrule
\end{tabular}
\end{center}
Bei gleicher Genauigkeit spart die gemeinsame Variante 10, 19 bzw.\ 16\,\% Speicher in den
Niedrigrangblöcken. Sie ist beim Aufbau schneller, und ihr Matrix-Vektor-Produkt
ebenfalls (0,14\,s gegenüber 0,20--0,24\,s). Das stützt den ersten Teil von H3 in mäßigem
Umfang. Die dichten Nahblöcke (29\,MB) sind bei beiden Varianten gleich.

\section{Skalierung}
Gemeinsame Variante, $\varepsilon=10^{-4}$, $k=1{,}5$:
\begin{center}\small
\begin{tabular}{rrrrrrr}
\toprule
$N$ & Unbekannte & $\mathcal{H}$ gesamt & davon Niedrigrang & dichter Kern & Aufbau & Mat-Vek\\
\midrule
720 & 5\,760 & 25,6\,MB & 4,7\,MB & 32\,MB & 17\,s & 0,07\,s\\
1280 & 10\,240 & 53,7\,MB & 25,1\,MB & 100\,MB & 37\,s & 0,14\,s\\
2000 & 16\,000 & 114,4\,MB & 37,8\,MB & 244\,MB & 55\,s & 0,25\,s\\
2880 & 23\,040 & 166,7\,MB & 89,9\,MB & 506\,MB & 93\,s & 0,35\,s\\
3920 & 31\,360 & 269,9\,MB & 120,9\,MB & 938\,MB & 124\,s & 0,53\,s\\
\bottomrule
\end{tabular}\\[2pt]
{\footnotesize Die dichte $E$-Matrix mit $8\times8$-Blöcken hätte bei $N=3920$ etwa 15\,GB.
Fehler in allen Fällen $0{,}85$ bis $1{,}1\cdot10^{-5}$.}
\end{center}
Der Kompressionsfaktor gegenüber dem dichten Kern steigt von $1{,}2$ auf $3{,}5$. Eine
Ausgleichsgerade ergibt Speicher $\propto N^{1{,}39}$ und Aufbauzeit $\propto N^{1{,}18}$;
$\text{Speicher}/(N\log N)$ steigt noch leicht an. Der Bereich ist also präasymptotisch:
Mit wachsendem $N$ werden immer größere Blöcke zulässig, und die strenge Bedingung
$\operatorname{dist}>3h$ hält viele Blätter dicht (bei $N=3920$ sind 55\,\% des Speichers
dichte Nahblöcke). Nahezu lineare Komplexität (H3, zweiter Teil) ist damit weder
gezeigt noch widerlegt. Nötig sind größere $N$ und eine Abstimmung von Blattgröße und
Zulässigkeit.

\section{Streuung jenseits der dichten Grenze}
Gelöst wird $T_1h=h^{\mathrm{inc}}$ mit GMRES (Neustart 200, Toleranz $10^{-6}$) und der
punktweisen Vorkonditionierung $2(1+J)^{-1}$. Beide Operatoren $E_{k_1}$ und $E_{k_2}$ sind
als $\mathcal{H}$-Matrizen gespeichert. Die Werte für $N\le500$ stammen aus der dichten
Rechnung in AP~2.
\begin{center}\small
\begin{tabular}{rrcccc}
\toprule
& & \multicolumn{2}{c}{Glas, $\omega a=1$} & \multicolumn{2}{c}{Gold, $\omega a=0{,}5$}\\
$N$ & Unbekannte & $Q_{\mathrm{ext}}$ & GMRES & $Q_{\mathrm{ext}}$ & GMRES\\
\midrule
500 (dicht) & 4\,000 & 0,20721 & 12 & 0,64732 & 35\\
720 & 5\,760 & 0,20959 & 10 & 0,63429 & 26\\
1280 & 10\,240 & 0,21198 & 10 & 0,61987 & 25\\
2000 & 16\,000 & 0,21310 & 10 & 0,61226 & 25\\
2880 & 23\,040 & 0,21371 & 10 & 0,60764 & 25\\
\midrule
\multicolumn{2}{l}{extrapoliert ($N$ = 1280, 2000, 2880)} & 0,21511 & & 0,59247 & \\
\multicolumn{2}{l}{Mie} & 0,21510 & & 0,59002 & \\
\bottomrule
\end{tabular}
\end{center}
Die GMRES-Zahlen der dichten Zeile bei $N=500$ sind ohne Neustart und mit Toleranz
$10^{-8}$ ermittelt und daher nicht direkt vergleichbar. Mit $\mathcal{H}$-Matrizen und
gleicher Einstellung sind die Iterationszahlen von $N$ unabhängig.

Für Glas ist die Konvergenz sauber zweiter Ordnung ($p=1{,}99$), und die Extrapolation
trifft Mie auf $4\cdot10^{-5}$. Für Gold ist die beobachtete Ordnung $p\approx1{,}46$
geringer, und die Extrapolation liegt $0{,}4\,\%$ über Mie. Die Ursache ist noch offen.
Plausibel sind die Kanten des Polyeders, auf die plasmonische Antworten empfindlich
reagieren, sowie die langsam konvergente Nahfeldquadratur (AP~2). Die Komprimierung ist
nicht die Ursache, denn ihr Fehler liegt bei $10^{-5}$.

\section{Würfel: Kompression und Kanten-/Eck-Vorkonditionierung kombiniert}\label{sec:kombi}

\paragraph{Netz.} Damit die Kantenblöcke aus AP~2.5 nicht entlang der Kante wachsen, wird
jede Seitenfläche in drei Zonen zerlegt (\texttt{cube\_aniso.py}): ein grobes Innenquadrat
($4\times4$), vier Kantenstreifen, die nur \emph{quer} zur Kante geometrisch gradiert sind
(Breite $0{,}5$, $L$ Schichten, entlang der Kante 4 Elemente), und vier Eckquadrate, die in
beiden Richtungen gradiert sind. An den Zonengrenzen entstehen hängende Knoten; für
stückweise konstante $L^2$-Ansätze ist das zulässig. Die Spurtrennung gilt auf dem
nichtkonformen Netz mit 2 bis 7\,\% Fehler bei 432 Dreiecken. Das ergibt 768, 1200, 1728
und 2352 Dreiecke für $L=2,\dots,5$.

\paragraph{Löser.} $E_{k_1}$ und $E_{k_2}$ werden als $\mathcal{H}$-Matrizen gespeichert
($\varepsilon=10^{-4}$, ohne Symmetriezerlegung). Die Kanten- und Eckblöcke (Radius
$R=0{,}25$) werden exakt mit Nahfeld assembliert und dicht per LU invertiert. Außerhalb der
Blöcke wird $\big(\tfrac12(1+J)\big)^{-1}$ verwendet. GMRES läuft rechtsvorkonditioniert,
mit Neustart 300 und Toleranz $10^{-6}$ (\texttt{cube\_h.py}), bei $\omega a=0{,}5$.
\begin{center}\small
\begin{tabular}{rrcccccc}
\toprule
& & \multicolumn{3}{c}{$\kappa=-3+0{,}3i$} & \multicolumn{3}{c}{Gold $-11+1{,}2i$}\\
$L$ & $N$ & ohne & Blöcke & $\sigma_{\mathrm{ext}}$ & ohne & Blöcke & $\sigma_{\mathrm{ext}}$\\
\midrule
2 & 768 & 73 & 49 & 18,235 & 64 & 45 & 19,909\\
3 & 1200 & 84 & 50 & 18,686 & 67 & 47 & 19,651\\
4 & 1728 & 96 & 49 & 18,593 & 69 & 47 & 19,404\\
5 & 2352 & 106 & 49 & 18,326 & 71 & 47 & 19,243\\
\bottomrule
\end{tabular}\\[2pt]
{\footnotesize $\mathcal{H}$-Speicher beider Operatoren: 50, 99, 169, 254\,MB. Größter Block:
192, 416, 608, 800 Unbekannte. Nur mit $2(1+J)^{-1}$: 66 bzw.\ 75 ($-3+0{,}3i$) und 57 bzw.\ 61
(Gold) für $L=2,3$. $\sigma_{\mathrm{ext}}$ mit und ohne Vorkonditionierung auf 6 Stellen
gleich.}
\end{center}

\paragraph{Befunde.}
\begin{itemize}[leftmargin=*]
  \item \emph{Iterationen:} Mit Kanten- und Eckblöcken sind sie konstant (49 bzw.\ 47), auch
    mit komprimierten Operatoren und über alle Symmetriesektoren zugleich. Ohne Blöcke
    wachsen sie um etwa 11 ($-3+0{,}3i$) bzw.\ 2--3 (Gold) je Stufe. Damit ist AP~2.5 mit AP~3
    kombiniert.
  \item \emph{Genauigkeit:} Für Gold fällt $\sigma_{\mathrm{ext}}$ monoton (Differenzen
    $0{,}26$; $0{,}25$; $0{,}16$). Für $-3+0{,}3i$ schwankt der Wert
    ($18{,}24\to18{,}69\to18{,}59\to18{,}33$). Das entspricht qualitativ der Vorhersage aus
    AP~1: Gold liegt außerhalb des Resonanzfensters, $-3+0{,}3i$ nahe daran. Da nur die
    Kantenzonen verfeinert werden (Innen- und Längsauflösung sind fest), konvergiert die
    Folge gegen die Lösung eines Netzes mit fester grober Restauflösung, nicht gegen die
    exakte Lösung. Ein Referenzwert für den Würfel fehlt.
  \item \emph{Blockanteil:} Bei $L=5$ decken die Blöcke $82\,\%$ der Unbekannten ab, weil in
    diesem Netz fast alle Elemente in den verfeinerten Zonen liegen. Die Blöcke selbst sind
    klein (höchstens 800 Unbekannte), die Vorkonditionierung ist hier aber nahezu ein
    blockweiser Direktlöser. Bei feinerem Innenbereich sinkt der Anteil. Eine
    $\mathcal{H}$-LU der Blöcke (AP~3.4 im engeren Sinn) war bei diesen Größen nicht nötig
    und ist nicht umgesetzt.
  \item \emph{Aufbauzeit:} Bei $L=5$ etwa 215\,s je Wellenzahl, mit 3216 dichten gegenüber
    2590 Niedrigrangblöcken. Ursache ist die Zulässigkeitsbedingung
    $\operatorname{dist}>3h_{\max}$: Die gestreckten Kantendreiecke haben lange Kanten (0,25)
    und machen viele Blöcke unzulässig. Nötig ist eine Nahfeldquadratur, die auch für
    gestreckte Elemente in kleinerem Abstand genau ist, oder eine Zulässigkeit nach der
    kurzen Elementausdehnung mit angepasster Fernfeldquadratur.
\end{itemize}

\section{Asymptotik mit dem C++-Rechenkern}\label{sec:cpp}
Der Kern aus dem Repository \texttt{PhotonPower/CliffordBem} (v0.1) implementiert dieselben
Kernkomponenten, Nahfeldregeln und ACA-Varianten. Die Einträge stimmen mit dem Prototyp auf
$10^{-13}$ überein, und er ist etwa 6-mal schneller. Einstellungen wie oben ($\varepsilon=10^{-4}$,
gemeinsame ACA). Details und Rohdaten: \texttt{docs/results\_compression.md} im Repository.
\begin{center}\small
\begin{tabular}{lrrrrr}
\toprule
Geometrie & Dreiecke & Unbekannte & Speicher & $10^3\cdot\mathrm{MB}/(N\log^2N)$ & Aufbau\\
\midrule
Kugel, $k=1{,}5$ & 1\,280 & 10\,240 & 55\,MB & 0,83 & 6\,s\\
 & 5\,120 & 40\,960 & 340\,MB & 0,91 & 33\,s\\
 & 11\,520 & 92\,160 & 959\,MB & 0,95 & 80\,s\\
 & 20\,480 & 163\,840 & 1\,881\,MB & 0,93 & 156\,s\\
Würfel (gleichm.), $k=1{,}5$ & 1\,200 & 9\,600 & 51\,MB & 0,84 & 7\,s\\
 & 6\,912 & 55\,296 & 488\,MB & 0,90 & 53\,s\\
 & 19\,200 & 153\,600 & 1\,651\,MB & 0,88 & 171\,s\\
Kugel, $k=n/4$ & 1\,280 & 10\,240 & 55\,MB & 0,83 & 6\,s\\
 & 20\,480 & 163\,840 & 2\,075\,MB & 1,03 & 192\,s\\
\bottomrule
\end{tabular}
\end{center}
\begin{itemize}[leftmargin=*]
  \item Der Quotient Speicher$/(N\log^2N)$ ist für Kugel und Würfel ab etwa 3000 Dreiecken
    innerhalb von $\pm7\,\%$ konstant. Die Aufbauzeit wächst mit Exponent $1{,}1$ bis
    $1{,}2$. Die frühere Beobachtung $N^{1{,}39}$ (Abschnitt~3) war präasymptotisch.
  \item Kanten und Ecken ändern das Wachstum nicht.
  \item Bei fester Auflösung je Wellenlänge (etwa 10 Elemente) kostet die Frequenz bis
    $kD\approx16$ nur etwa $10\,\%$ Speicher. Bei festem Netz wächst der Niedrigrang-Speicher
    bis $kD=48$ um $62\,\%$, und der mittlere Rang steigt von 7,0 auf 9,0.
  \item Gemeinsame gegen komponentenweise ACA: etwa $19\,\%$ weniger
    Niedrigrang-Speicher, konstant über $N=1280$ bis $11\,520$, bei halber Aufbauzeit.
  \item Zu den Kanten gradierte Tensornetze sind dicht-dominiert (bei 3072 Dreiecken $84\,\%$
    dichte Nahblöcke, 293\,MB gegenüber 170\,MB gleichmäßig). Das liegt an den auf $h_{\max}$
    bezogenen Nahfeld- und Zulässigkeitskriterien und ist mit einer Fernfeldquadratur für
    gestreckte Elemente zu beheben.
\end{itemize}
\paragraph{Streulöser.} Der C++-Kern (v0.2) löst $T_1h=h^{\mathrm{inc}}$ mit beiden Operatoren als
$\mathcal{H}$-Matrizen und GMRES (\texttt{docs/results\_scattering.md}). Bei 320 Dreiecken
reproduziert er die dichte Lösung des Prototyps auf alle Stellen. Bis 92\,160 Unbekannte:
\begin{center}\small
\begin{tabular}{lcccc}
\toprule
& $Q_{\mathrm{ext}}$ (11\,520 Dreiecke) & extrapoliert & Mie & GMRES\\
\midrule
Glas, $\omega a=1$ & 0,214752 & 0,215111 ($p=1{,}96$) & 0,215098 & 10\\
Gold, $\omega a=0{,}5$ & 0,597669 & 0,590837 ($p=1{,}29$) & 0,590018 & 24\\
\bottomrule
\end{tabular}
\end{center}
Die Iterationszahlen sind von $N$ unabhängig. Bei Gold sinkt die beobachtete Ordnung mit
$N$ ($1{,}42\to1{,}29$), die Extrapolation nähert sich Mie ($3{,}2\to1{,}4\cdot10^{-3}$).

\paragraph{Sauter-Schwab (C++-Kern v0.4).} Mit Sauter-Schwab-Quadratur für alle Paare mit gemeinsamer
Fläche, Kante oder Ecke ändert sich das Bild für Gold grundlegend. Für ein kantenbenachbartes
Paar hatte die bisherige Nahfeldregel 1,7\,\% Fehler, Sauter-Schwab (Ordnung 5) liegt unter
$10^{-5}$. Die Streurechnung konvergiert für Gold jetzt mit Ordnung $p=1{,}99$--$2{,}02$, und die
Extrapolation trifft Mie auf $2\cdot10^{-5}$ bis $9\cdot10^{-5}$ (zuvor $1{,}4\cdot10^{-3}$). Auf dem
feinsten Netz sinkt der Fehler ohne Extrapolation von 1,3\,\% auf 0,19\,\%. Die langsamere
Konvergenz bei Gold lag also an der Nahfeldquadratur, nicht an Formulierung oder Ansatzraum.

Damit ist H3 (zweiter Teil) im Bereich bis etwa 160\,000 Unbekannte und $kD\lesssim16$ numerisch
gestützt: Der Aufwand ist $\mathcal{O}(N\log^\alpha N)$ mit $\alpha\approx2$ für den Speicher.

\section{Weiterentwicklung des C++-Kerns (v0.4 bis v0.22)}\label{sec:cpp2}
Kurzfassung; Einzelheiten in den Berichten unter \texttt{docs/} im Repository.
\begin{itemize}[leftmargin=*]
  \item \emph{Nahfeld.} Sauter-Schwab-Quadratur für benachbarte Paare: Die bisherige Regel hatte an gemeinsamen
    Kanten bis 1,7\,\% Fehler; mit Sauter-Schwab konvergiert Gold mit Ordnung 2 gegen Mie (Extrapolation
    $2\cdot10^{-5}$ bis $9\cdot10^{-5}$ statt $1{,}4\cdot10^{-3}$). Auf gestreckten Dreiecken (Seitenverhältnis bis 16)
    ist Sauter-Schwab ungenau (bis 1,4\,\%) und orientierungsabhängig; dort wird eine halbanalytische, zur
    gemeinsamen Kante gradierte Regel verwendet, für nahe Paare eine adaptive Außenregel, alles in einem Nahfeld-Cache.
  \item \emph{Chirale Medien.} Innerer Operator $P_+E_{k_+}+P_-E_{k_-}$, chirale $J$; $Q_\pm$ und CD treffen eine
    chirale Mie-Lösung nach Extrapolation auf $10^{-4}$.
  \item \emph{Mehrere Körper, Gmsh, Spektren.} Blockdiagonaler Innenoperator je Körper; Gmsh-Import;
    Johnson-Christy-Daten, Orientierungsmittelung (Lebedev). Born-Kuhn-Dimer: bisignates CD-Spektrum,
    Enantiomere mit entgegengesetztem CD auf $10^{-5}$.
  \item \emph{Vorkonditionierung.} Blockvorkonditionierung für mehrere Körper und chirale Medien (voller Block =
    1 Iteration); Würfel-Dimer mit Kanten-/Eckblöcken: Iterationen nahezu unabhängig von der Kantenauflösung
    ($56\to54$ statt $56\to72$). Hierarchische Faktorisierung im HODLR-Format (schwache Zulässigkeit): direkter
    Löser mit feiner Toleranz, als Vorkonditionierer lohnend ab etwa 7 bis 18 rechten Seiten; an Plasmonresonanzen
    nur begrenzt wirksam.
  \item \emph{ACA mit Multivektor-Pivots.} Kreuzapproximation über $\clc$ mit Nullteiler-Erkennung: bei gleicher
    Genauigkeit 2,3- bis 2,9-mal mehr Speicher als die gemeinsame skalare ACA (Kugel, Würfel, $k=12$). Der Kern ist
    skalar separabel (skalare Funktionen von $x$ mal paravektorwertige von $y$); genau das nutzt die gemeinsame ACA.
  \item \emph{Abgerundete Würfel.} Im Resonanzfenster konvergiert die Rechnung mit gerundeten Kanten monoton, hängt
    aber stark vom Radius ab; Silberwürfel (50\,nm, Wasser): Hauptresonanz verschiebt sich mit schärferen Kanten nach
    Rot (grob 440, 460, 470\,nm für 10, 7,5, 5\,nm Radius).
  \item \emph{Beschichtete Grenzflächen} (v0.13). Verschachtelte Gebiete über einen Grenzflächengraph: jede
    Schichtgrenze ist eine Fläche, jedes Gebiet hat einen eigenen komprimierten Cauchy-Operator auf seinem Rand
    (AP~1, Nachtrag). Beschichtete Kugel gegen Aden-Kerker: Extinktion und komplexe Vorwärtsamplitude mit Ordnung~2.
    Für Schichten dünner als die Elemente wird die Schichtwirkung als Differenz zu einer neutralen Rechnung auf
    denselben Netzen bestimmt (Fehler 1--5\,\% bis $d/h\approx0{,}04$, Phase 3--9\,\% auf dem feinsten Netz).
    Parallele Flächen im Abstand $d\ll h$: Die adaptive Außenregel verfeinert nur noch zu den Rändern des inneren
    Dreiecks hin (je Paar Faktor 3--4 schneller bei gleicher Genauigkeit, im Gesamtproblem etwa 20\,\% weniger
    Nahfeldzeit). Silberwürfel (50\,nm) mit 2\,nm Oxid: Rotverschiebung der Hauptresonanz um etwa 20\,nm,
    doppelt so viel wie bei der Kugel.
  \item \emph{Dünnschicht-Näherung} (v0.14). Eine Fläche mit $J_{\mathrm{eff}}=J+L_d$ (Greensche Funktion der
    Schichtfolge in erster Ordnung, Ableitungen auf der Dichte): Die $\mathcal{H}$-Matrizen der unbeschichteten Rechnung
    werden unverändert verwendet, $L_d$ ist ein lokaler Operator über die Kantennachbarn. Kosten wie ohne Schicht;
    Modellfehler $\approx1{,}5\,d/a$; für die Phase bei $d/a\lesssim0{,}03$ genauer als die Zwei-Flächen-Rechnung, bei etwa einem Fünftel der Kosten.
  \item \emph{Dünnschicht zweiter Ordnung} (v0.15). Dirac-Form auf Parallelflächen mit Formoperator; die zweiten
    Ableitungen sind zusammengesetzte Kleinste-Quadrate-Gradienten, alle Schichtkorrekturen nutzen glatte Normalen.
    Schichtwirkung bis $d/a=0{,}05$ auf 0,1--0,7\,\%, Silberkugel mit 2\,nm Oxid so genau wie die exakte Rechnung. Die
    zusätzlichen Kosten sind lokale Operatoren; die Iterationen steigen für $d\gtrsim h$.
  \item \emph{Konsistente zweite Ableitungen, mehrere Körper} (v0.16). Quadratische Anpassung über die
    Knotennachbarschaft mit geschlossener Form für $B^2$; blockdiagonaler Innenoperator für mehrere Körper. Am
    Silberwürfel braucht die Näherung eine aufgelöste Krümmung (Unterschied zur exakten Rechnung fällt von 11 auf 6\,\%
    bei Verfeinerung von 1\,208 auf 2\,048 Dreiecke).
  \item \emph{Chirale Schichten} (v0.17). Mit dem zentralen Multivektor $K=k_+P_++k_-P_-$ gilt die Dirac-Form
    unverändert; der CD einer dünnen chiralen Hülle auf Gold konvergiert mit Ordnung~2 gegen eine Mie-Lösung für
    geschichtete chirale Kugeln (0,5--1\,\% bei $n=12$), die exakte Zwei-Flächen-Rechnung verfehlt ihn um 30--50\,\%.
    Spektrum einer Goldkugel mit chiraler Molekülhülle (v0.18): plasmoninduzierter CD an der Resonanz auf 2--7\,\%.
  \item \emph{Schicht als Zweitor} (v0.19, Prototyp). Nach dem Vorbild der S-Matrix von Schichtsystemen: zwei
    $\mathcal{H}$-Matrizen ($E_2$ auf der Außenfläche, $E_1$ auf dem Kern) und eine Zweitor-Beziehung mit
    $g(s)=\tanh(d\sqrt s/2)/\sqrt s$, ausgewertet über Partialbrüche mit dünnbesetzten Resolventen des
    Laplace-Beltrami-Operators (5--11 innere Schritte je Pol). Stabil bis $d/h=3{,}8$ bei unter 62 Iterationen.
  \item \emph{Mehrfachschichten im Zweitor} (v0.20). Blockbidiagonales System mit den Spuren auf den inneren
    Grenzflächen; weiterhin nur zwei $\mathcal{H}$-Matrizen. Jede Schicht auf ihrer eigenen Fläche, dicke Schichten
    automatisch unterteilt: konvergent bis $d/a=0{,}5$, die Iterationen wachsen mit der Zahl der Teilschichten (bis 156).
  \item \emph{Chirales Außenmedium} (v0.21). Der Außenoperator wird $P_+E_{k_+}+P_-E_{k_-}$ (doppelte Zahl der
    Außen-$\mathcal{H}$-Matrizen), die Anregung eine Helizitätswelle, die Extinktion das optische Theorem je Kanal. Die
    Blockvorkonditionierung ist dafür noch nicht angepasst (HODLR und punktweise funktionieren).
  \item \emph{Mehrere Körper im Zweitor} (v0.22): $E_2$ auf der Vereinigung der Außenflächen, $E_1$ blockdiagonal; der
    Aufwand bleibt bei zwei (chiral vier) Operatoren. Die $\mathcal{H}$-Matrizen der Schichtgebiete
    enthalten mehr unzulässige Blöcke, weil sich die Cluster beider Flächen überlappen; Speicher und Nahpaare etwa
    verdoppeln sich je Schicht.
\end{itemize}

\section{Einordnung und nächste Schritte}
\begin{itemize}[leftmargin=*]
  \item \emph{H3, erster Teil:} Gemeinsame Pivotisierung über die Grade spart
    10--19\,\%. Multivektor-Pivots verbessern das nicht, sie kosten 2,3- bis 2,9-mal mehr Speicher
    (Abschnitt~\ref{sec:cpp2}); in dieser Form ist die Hypothese widerlegt.
  \item \emph{H3, zweiter Teil:} Mit dem C++-Kern bis 163\,840 Unbekannte gestützt:
    $\mathcal{O}(N\log^2N)$ für Kugel und Würfel (Abschnitt~\ref{sec:cpp}). Offen sind stark
    gradierte Netze (Kriterien für gestreckte Elemente) und elektrisch große Streuer
    ($kD\gg50$).
  \item \emph{Laufzeit:} Der C++-Kern ist etwa 6-mal schneller als der Python-Prototyp
    (ein Rechenkern); die Assemblierung ist für OpenMP vorbereitet.
  \item \emph{AP~3.4:} Kompression und Kanten-/Eck-Vorkonditionierung sind am Würfel
    kombiniert (Abschnitt~\ref{sec:kombi}). Offen sind die $\mathcal{H}$-LU für große Blöcke und
    eine Zulässigkeit, die gestreckte Kantendreiecke nicht bestraft.
\end{itemize}

\end{document}
